\section{Experiments}
\label{sec:experiments}

\subsection{Dataset and label semantics}

The dataset used in this work contains 27 CSV files and 37,417 timestamps sampled at a fixed interval of 4 seconds. Each row contains multi-sensor measurements and three label-related fields, among which \texttt{running\_state\_five\_class} is the main target. The verified state counts are:
\begin{itemize}
\item stop (`1`): 10,959,
\item normal (`5`): 15,695,
\item first-level degradation (`7`): 5,364,
\item second-level degradation (`9`): 185,
\item fault (`3`): 5,214.
\end{itemize}

The binary fault indicator \texttt{JianSuDuan\_ChaoSu} equals 1 only when the final fault state (`3`) is active, which supports the interpretation that the remaining labels are pre-fault or non-fault operational states within the same overspeed lifecycle.

\subsection{Evaluation protocol}

The main evaluation must use file-level generalization:
\begin{itemize}
\item leave-one-file-out cross-validation as the primary protocol,
\item grouped K-fold by file as a supplementary protocol.
\end{itemize}

Random window split is not used as the main setting because it overestimates generalization under strong temporal correlation.

Two tasks are evaluated:
\begin{enumerate}
\item current five-state recognition,
\item short-horizon warning with horizons of 12 s, 20 s, and 40 s.
\end{enumerate}

For state recognition, the main metrics are macro-F1, balanced accuracy, and per-class F1. For short-horizon warning, the main metrics are AUROC, AUPRC, event-level recall, false alarm rate, and time-bucket accuracy.

\subsection{Baselines and ablations}

The benchmark suite should include:
\begin{itemize}
\item flat 5-class TCN,
\item flat 5-class Transformer,
\item class-balanced TCN,
\item class-balanced Transformer,
\item SGPH-Net.
\end{itemize}

Mandatory ablations include:
\begin{itemize}
\item SGPH-Net without the state-graph constraint,
\item SGPH-Net without the hazard head,
\item SGPH-Net without the time-bucket head,
\item SGPH-Net without sensor grouping.
\end{itemize}

\subsection{Result tables to be filled after benchmark runs}

Table~\ref{tab:main_results} is reserved for the main comparison. The cells are intentionally left blank in this draft to avoid fabricating results before controlled runs.

\begin{table}[t]
\centering
\caption{Main comparison on five-state recognition.}
\label{tab:main_results}
\begin{tabular}{lccc}
\toprule
Method & Macro-F1 & Balanced Acc. & Fault F1 \\
\midrule
Flat TCN & -- & -- & -- \\
Flat Transformer & -- & -- & -- \\
Balanced TCN & -- & -- & -- \\
Balanced Transformer & -- & -- & -- \\
SGPH-Net & -- & -- & -- \\
\bottomrule
\end{tabular}
\end{table}

\begin{table}[t]
\centering
\caption{Short-horizon warning results.}
\label{tab:warning_results}
\begin{tabular}{lccc}
\toprule
Method & AUROC & AUPRC & Time-Bucket Acc. \\
\midrule
Balanced TCN & -- & -- & -- \\
Balanced Transformer & -- & -- & -- \\
SGPH-Net & -- & -- & -- \\
\bottomrule
\end{tabular}
\end{table}

\subsection{What the final paper must demonstrate}

The final version of the paper should demonstrate three empirical points:
\begin{enumerate}
\item SGPH-Net outperforms flat classifiers on macro-F1 or balanced accuracy.
\item SGPH-Net improves recognition of degradation states, especially the rare second-level degradation state.
\item SGPH-Net provides earlier and more stable warning than baselines at short horizons.
\end{enumerate}

Without these three points, the manuscript remains a method proposal rather than a submission-ready conference paper.
